\documentclass[12pt]{article}
\usepackage[T1]{fontenc}
\usepackage[utf8]{inputenc}
\usepackage{lmodern}
\usepackage{textcomp}
\usepackage{enumitem}
\usepackage{geometry}
\geometry{margin=1in}
\begin{document}
\begin{center}
Answer Key - Political Science - Graduate Level Assessment 6 \
\small (Answers are ordered exactly as in the question paper.)
\end{center}

\section*{Multiple Choice}
\begin{enumerate}[label=\arabic*.]
\item \textbf{Answer:} D
\\ \textit{Options:} A) Limiting security analysis to military engagement is too restrictive.; B) Security is a contestable concept.; C) Focusing security only on violent conflicts is too restrictive.; D) All of these options.
\\ \textit{Note:} The correct answer is "All of these options" because the Human Development Report has effectively integrated various dimensions of human security-such as economic, political, and social factors-into a comprehensive framework that emphasizes the interconnectedness and holistic approach necessary for addressing global security challenges.
\item \textbf{Answer:} D
\\ \textit{Options:} A) Any crime that uses computers to jeopardise or attempt to jeopardise national security; B) The use of computer networks to commit financial or identity fraud; C) The theft of digital information; D) Any crime that involves computers and networks
\\ \textit{Note:} The term 'cyber-crime' refers to illegal activities that are facilitated or committed through the use of computers and networks, highlighting the integral role of technology in modern criminal behavior.
\item \textbf{Answer:} D
\\ \textit{Options:} A) The International Monetary Fund (IMF); B) The US Dollar as the world's reserve currency; C) The General Agreement on Tariffs and Trade (GATT); D) All of the above
\\ \textit{Note:} The correct answer "All of the above" is justified because the Bretton Woods Conference in 1944 led to the establishment of key institutions such as the International Monetary Fund (IMF) and the World Bank, which collectively aimed to promote international financial stability and economic cooperation.
\item \textbf{Answer:} C
\\ \textit{Options:} A) Creation of the Department of Homeland Security; B) Moving the Department of Defense Intelligence Agencies to the CIA for better coordination; C) Creation of the National Intelligence Director; D) Moving the domestic intelligence component of the FBI to the CIA
\\ \textit{Note:} The 9/11 Commission proposed the creation of the National Intelligence Director to centralize and enhance the coordination of the U.S. intelligence community, addressing the fragmented oversight that contributed to the failures leading up to the September 11 attacks.
\item \textbf{Answer:} D
\\ \textit{Options:} A) United States, Soviet Union, Germany, France, and Great Britain; B) United States, Germany, France, Great Britain, and Japan; C) United States, Great Britain, Republic of China, India, and Brazil; D) United States, Soviet Union, France, Great Britain, and Republic of China
\\ \textit{Note:} The original five permanent members of the UN Security Council-United States, Soviet Union, France, Great Britain, and Republic of China-were designated as such in the UN Charter in 1945 due to their significant roles in World War II and their status as major global powers at the time, granting them veto power over substantive resolutions.
\end{enumerate}

\section*{True / False}
\begin{enumerate}[label=\arabic*.]
\setcounter{enumi}{5}
\item \textbf{Answer:} True
\\ \textit{Options:} A) True; B) False
\\ \textit{Note:} The statement is true because GATT (General Agreement on Tariffs and Trade), the IMF (International Monetary Fund), and the World Bank are foundational institutions that collectively promote international trade, financial stability, and economic development, which are central to the principles of the liberal international economic order.
\item \textbf{Answer:} False
\\ \textit{Options:} A) True; B) False
\\ \textit{Note:} The Roosevelt Doctrine, articulated in 1904, asserts that the United States will intervene in Latin American countries to maintain stability and prevent European intervention, thereby discouraging rather than encouraging European involvement in the region.
\item \textbf{Answer:} False
\\ \textit{Options:} A) True; B) False
\\ \textit{Note:} Congress opposed Wilson's proposal for the League of Nations primarily due to concerns over national sovereignty and the potential entanglement in international conflicts, rather than a direct fear of undermining U.S. economic interests.
\item \textbf{Answer:} False
\\ \textit{Options:} A) True; B) False
\\ \textit{Note:} The 'Philadelphian System' primarily refers to the framework established by the U.S. Constitution that emphasized domestic governance and federalism, which limited direct involvement in European political affairs and promoted neutrality instead.
\item \textbf{Answer:} True
\\ \textit{Options:} A) True; B) False
\\ \textit{Note:} Nations are shaped by shared cultural identities, historical experiences, and social bonds that transcend geographical boundaries, highlighting that the essence of a nation lies in its people and their collective consciousness rather than merely in political delineations.
\end{enumerate}

\section*{Long Form Response}
\begin{enumerate}[label=\arabic*.]
\setcounter{enumi}{10}
\item \textbf{Answer:} Postcolonialism identifies three primary forms of violence: physical violence, structural violence, and cultural violence. Physical violence refers to direct acts of harm inflicted on individuals or groups, while structural violence encompasses systemic inequalities that perpetuate suffering and disadvantage among marginalized populations. Cultural violence involves the ways in which dominant groups promote ideologies that justify and normalize oppression. Pervasive violence, although omnipresent in postcolonial contexts, is not categorized as a distinct form because it lacks the specificity and agency found in the other three; rather, it is an overarching condition that results from the interplay of the identified forms, making it more of a backdrop than a separate category of analysis.
\\ \textit{Note:} correct because it accurately defines the three forms of violence in postcolonial theory-physical, structural, and cultural-while also clarifying that pervasive violence is not included as a distinct category due to its omnipresence rather than its specific mechanisms of oppression and harm.
\item \textbf{Answer:} Offshore Balancing is characterized by several key policies that significantly shape international relations and state security responsibilities. Central to this strategy is the embrace of multi-polarity, which encourages the participation of various regional powers in maintaining stability, rather than relying solely on a dominant state. This approach advocates for greater restraint in international interventions, allowing the United States to prioritize its resources and focus on strategic interests rather than entangling itself in every regional conflict. Additionally, Offshore Balancing necessitates that other states take on their own security burdens, promoting a sense of accountability and encouraging them to develop their military capabilities. Overall, these policies reflect a pragmatic shift towards a more sustainable and collaborative international order, where the burden of security is equitably shared among nations.
\\ \textit{Note:} correct because it identifies the central tenets of Offshore Balancing, such as the promotion of multi-polarity and strategic restraint in interventions, which collectively influence the dynamics of international relations and the allocation of state security responsibilities.
\end{enumerate}

\vfill
\noindent\textit{End of Answer Key}
\end{document}